% \documentclass{article}
% \usepackage[utf8]{inputenc}
% \usepackage{amsmath}
% \usepackage{graphicx}
% \usepackage{booktabs}
% \usepackage{array}
% \usepackage{geometry}\geometry{a4paper, margin=1in}
% \usepackage{hyperref}
% \usepackage{parskip}

% \title{Architectural Description: Human Benchmark \\ \Large Context View}
% \author{Generated by AI Software Architect Assistant}
% \date{\today}

% \begin{document}

% \maketitle
% \tableofcontents

\section{Introduction to the Context View}

The Context View for the Human Benchmark system defines
its environment, establishes its boundaries
and identifies all external entities with which the system
interacts. It provides essential insight about system’s scope, 
helping stakeholders understand exactly what is included and what lies beyond it.

\subsection{Primary Stakeholders}

The primary stakeholders for this Context View are:

\begin{itemize}
    \item \textbf{Acquirers}: Those funding and approving decisions about the system's development, who need to understand system scope and external dependencies.
    \item \textbf{Users}: Players, game creators who interact with the system.
    \item \textbf{Developers}: Who need to understand the system, its boundaries and external interfaces.
    \item \textbf{System Administrators}: Responsible for deploying and maintaining the system and its connections.
    \item \textbf{Testers}: Who need to understand system boundaries for comprehensive testing.
\end{itemize}

All stakeholders benefit from this view as it establishes a common understanding of what the system is and is not, along with its interactions with the outside world.

\section{Core Architectural Elements of the Context View}

\subsection{System Scope and Responsibilities}

The Human Benchmark system is a web application designed to measure and enhance users' 
cognitive skills through various types of games. 
The application operates as both a training platform and a competitive environment where users can track their performance, compete with other users, friends and improve their reaction time, reflexes while also having fun.

\subsubsection{Key Capabilities}

The Human Benchmark system provides the following key capabilities:

\begin{itemize}
    \item User registration and authentication with secure password management
    \item Personal user profiles with performance tracking and game history
    \item Single-player cognitive and reaction tests/games with precise timing measurement
    \item Customizable game creation with adjustable parameters
    \item Multiplayer functionality with real-time competition
    \item Public and private game rooms
    \item Real-time communication
    \item Comprehensive score recording and leaderboard systems
    \item Game session tracking and management
\end{itemize}

\subsubsection{Exclusions}

The following capabilities are explicitly excluded from the current scope of the Human Benchmark system:

% LGTM

\begin{itemize}
    % ! \item Integration with external AI training systems
    \item Revenue generation or payment processing functionality (the system is
          currently in a bootstrapping phase, where initial operations are funded by
          saved capital or investment rather than immediate profits) % ! or payment
          processing functionality
    \item Third-party authentication providers (currently using email/password
          authentication)
    \item Mobile application versions (current focus is web-based)
          % ! \item Offline functionality (system requires network connectivity)
          % ! \item Integration with external health or cognitive assessment platforms
\end{itemize}

\subsection{External Entities and Services}

\subsubsection{Users and Roles}

\begin{itemize}
    \item \textbf{Players} - Primary users who access the system to play games,
          track performance, and compete with others. They interact with the system
          through web browsers.

          % ! \item \textbf{Administrators} - System operators who manage user
          % ! accounts, monitor system usage, and maintain operational status. They
          % ! interact with the system through dedicated administrative interfaces.

    \item \textbf{Game Creators} - A subset of users who create custom game
          configurations for themselves or other users. They need interfaces to define
          game parameters and game access settings.
\end{itemize}

\subsubsection{External Systems}

% LGTM

\begin{itemize}
    \item \textbf{Client Devices} - Various hardware devices (desktops, laptops,
          tablets) that users use to access the system. These vary in screen size,
          processing power affecting the user experience.

    \item \textbf{Web Browsers} - Different browsers through which users access
          the system. Browser compatibility and capabilities influence the system's
          presentation layer.

    \item \textbf{Database System (PostgreSQL)} - External database system for
          persistent storage of user data, game configurations, scores, and session
          information.

    \item \textbf{CI/CD Pipeline} - External build and deployment systems
          (GitLab Actions) that automate the building, testing, versioning, database
          updating and deployment of new system versions. The pipeline is executed by
          runner on MIF OpenNebula virtual machine.

    \item \textbf{Containerization Platform (Docker)} - External platform for
          packaging and distributing the system components. (In our case, it is done
          on the MIF OpenNebula virtual machine)

          % ! \item \textbf{Container Orchestration (Kubernetes)} - External system
          % ! anaging deployment, scaling, and operations of application containers.

    \item \textbf{Hosting Infrastructure (MIF OpenNebula)} - Virtual machine
          where CI/CD runner executes, the docker containers are deployed and the system
          is hosted.

\end{itemize}

% TODO: DIAGRAM - High-level context model showing Human Benchmark system and all its external entities
% SUGGESTED_TYPE: Custom block diagram (informal boxes-and-lines notation)
% KEY_ELEMENTS_TO_DEPICT: Human Benchmark System (center), Users (Players, Administrators), Client Devices, Web Browsers, Database System, CI/CD Pipeline, Hosting Infrastructure
\begin{figure}[h!]
    \centering
    \includegraphics[width=0.6\textwidth]{AutomationAndSections/Sections/1contextView/Images/systemContext.png}
    \caption{Human Benchmark System Entity Overview}
    \label{fig:context_system_overview}
\end{figure}

\subsection{Nature and Characteristics of External Entities}

\subsubsection{Client Devices and Browsers}

% The quality properties of external entities, such as system stability and
% availability, performance and throughput capabilities, physical location, or
% data quality, may significantly affect the architecture of the system.

% LGTM

\begin{itemize}
    \item \textbf{Availability:} Variable - devices can disconnect at any time
    \item \textbf{Performance:} Highly variable - ranging from low-end to high-end devices
    \item \textbf{Data Quality:} Variable - dependent on user input quality and browser capabilities
    \item \textbf{Location:} Geographically distributed - potentially worldwide
\end{itemize}

These properties influence the system's design to ensure responsive UI across
different devices and good handling of connectivity issues. The system must be
well optimized for content delivery accros various network conditions and
devices.

\subsubsection{Database System}

% LGTM

\begin{itemize}
    \item \textbf{Availability:} High - critical for system operation
    \item \textbf{Stability:} High - should be resilient to failures
    \item \textbf{Throughput capabilities:} Medium - moderate read/write operations
    %\item \textbf{Data Quality:} High - enforced through database constraints
    \item \textbf{Location:} Different from the application server - hosted on a cloud provider
\end{itemize}

The database system provides essential data services. Its
reliability directly impacts the system's overall availability, and
performance during peak usage.

\subsubsection{CI/CD Pipeline and Infrastructure Components}

% LGTM

\begin{itemize}
    \item \textbf{Availability:} Medium - should be available for deployment and testing
    \item \textbf{Performance:} Medium - application and CI/CD pipeline runs on one virtual machine. These must not interfere with each other.
    \item \textbf{Location:} Same as the application server - hosted on a MIF OpenNebula virtual machine
\end{itemize}

These systems enable rapid and reliable deployment of application updates but are not directly involved in runtime operations.

\subsubsection{Hosting Infrastructure}
% LGTM

\begin{itemize}
    \item \textbf{Availability:} High - critical for game sessions and user interactions
    \item \textbf{Performance:} Medium - should support expected load
    \item \textbf{Scalability:} Medium - should support +200 concurrent users
    \item \textbf{Security:} High - must protect sensitive user data and system (database) login credentials must be stored securely
    \item \textbf{Location:} hosted on a MIF OpenNebula virtual machine
\end{itemize}

The hosting infrastructure is critical for the user experience.



\subsection{External Interfaces}

\begin{table}[h!]
\centering
\begin{tabularx}{\textwidth}{|l|X|X|X|}
\hline
\textbf{Interface} & \textbf{Type} & \textbf{Purpose} & \textbf{Technology} \\
\hline
\textbf{User} & Service provider (interactive web UI) & Allows users to register, authenticate, configure games, play games, view scores, and manage profiles & HTTPS, WebSockets for real-time communication \\
\hline
\textbf{Database} & Data provider and consumer & Persists and retrieves user data, game configurations, scores, and session information & Entity Framework Core over PostgreSQL \\
\hline
\textbf{WebSocket} & Service provider and consumer, Event provider and consumer & Provides real-time communication for multiplayer functionality, game state synchronization, and live updates & WebSocket protocol over HTTPS \\
\hline
\textbf{Deployment} & Service consumer & Handles application deployment and updates & Docker containers, GitLab CI/CD, MIF OpenNebula \\
\hline
\end{tabularx}
\caption{External interfaces for Human Benchmark System}
\label{tab:external_interfaces}
\end{table}



% Identity and Responsibilities of External Interfaces
% For each external entity, the nature of all interfaces between it and this sys-
% tem should be identified. Such an interface may serve one of the following
% purposes.
% Data provider or consumer: The external system supplies data directly to
% this system or receives data directly from it.
% Service provider or consumer: The external system is requested to per-
% form some action by this system or requests some action of this system
% (e.g., a service call), and the service may return data and/or status infor-
% mation in response to the request.
% Event provider or consumer: The external system publishes events that
% this system wishes to be notified of, or this system publishes events that
% the external system wishes to be notified of.

% * CHECK PAGE 251 for more details on the interface types



% Seemed pretty useless, so I commented it out

% !!! do not trust this blindly, check the original document
% TODO: DIAGRAM - Sequence diagram showing interaction between user and system for a typical game session
% SUGGESTED_TYPE: UML Sequence Diagram
% KEY_ELEMENTS_TO_DEPICT: User, Web Client, Human Benchmark Backend, Database, including authentication, game selection, game play, and score recording
%\begin{figure}[h!]
%    \centering
    % \includegraphics[width=0.9\textwidth]{images/placeholder_diagram_for_context_view_user_interaction.png}
%    \caption{Illustrative Diagram: User Interaction Sequence for Game Session (Conceptual Placeholder)}
%    \label{fig:context_user_interaction}
%    \textit{Comment: This diagram should illustrate the sequence of interactions
%        between a user and the Human Benchmark system throughout a typical gaming
%        session.}
%\end{figure}






\subsection{Nature and Characteristics of External Interfaces}

\subsubsection{User Interface Characteristics}

% LGTM

\begin{itemize}
    \item \textbf{Characteristics:}
          \begin{itemize}
              \item Highly interactive with real-time feedback
                    % ! \item Responsive design for various screen sizes
              \item Varying latency requirements (extremely low for reaction
                    time tests)
                    % !  \item Automated timing and score calculation
          \end{itemize}
    \item \textbf{Volumes:}
          \begin{itemize}
              \item Peak concurrent users: 200
              \item Data size: Generally small payloads for game interactions
              \item Growth: Will not grow significantly in the near future
          \end{itemize}
    \item \textbf{Interaction Triggers:} Event-driven (user actions)
    \item \textbf{Automation Level:} Semi-automated (user initiates, system processes)
    \item \textbf{Criticality:} High for reaction time tests (millisecond precision required)
\end{itemize}

\subsubsection{Database Interface Characteristics}

% LGTM

\begin{itemize}
    \item \textbf{Characteristics:}
          \begin{itemize}
              \item Moderate volume of read operations (score queries)
              \item Moderate volume of write operations (new scores, user registrations)
              \item High reliability requirements (data integrity critical)
              \item Transactional requirements for score recording and user management
              \item User password hashing and encrypting for security
          \end{itemize}
    \item \textbf{Volumes:}
          \begin{itemize}
              \item User data: A few hundred records per month
              \item Game sessions: Potentially thousands of records per month
              \item Score data: Potentially a few thousand records per month
          \end{itemize}
    \item \textbf{Interaction Triggers:} Event-driven (user actions, system events)
    \item \textbf{Automation Level:} Fully automated
    \item \textbf{Transactionality:} High (critical for user data integrity)
\end{itemize}

\subsubsection{WebSocket Interface Characteristics}

% LGTM

\begin{itemize}
    \item \textbf{Characteristics:}
          \begin{itemize}
              \item Bidirectional communication
              \item Event-driven interaction model
              \item Varying connection stability (dependent on user network conditions)
              \item Scaling challenges with high concurrent user counts
              \item Important for real-time game state updates
          \end{itemize}
    \item \textbf{Volumes:}
          \begin{itemize}
              \item Connection count: Equal to active game session user count
              \item Message frequency: High during active games (multiple per second)
              \item Message size: Small (only information about target changes)
          \end{itemize}
    \item \textbf{Interaction Triggers:} Real-time events (user clicks on targets)
    \item \textbf{Automation Level:} Semi-automated (user initiates, system processes)
    \item \textbf{Criticality:} High for competitive gameplay (low latency required)
\end{itemize}

\subsubsection{Deployment Interface Characteristics}

% LGTM

\begin{itemize}
    \item \textbf{Characteristics:}
          \begin{itemize}
              \item Scheduled deployments (idealy durring off-peak hours)
                    %   \item Automated deployment processes
                    %   !\item Rolling updates to minimize downtime
                    %   \item Containerized environment for consistency
                    %   !\item Automated scaling based on demand
                    %   \item Automated building, testing, and deployment pipeline processes
                    %   \item Manual versioning and database update triggering pipeline processes
                    %   \item Only authorized personnel can trigger deployments
              \item Transactional - all of the stages must be completed successfully before the deployment
                    is considered successful
              \item Database login credentials and https certificates must be passed securely
          \end{itemize}
    \item \textbf{Volumes:}
          \begin{itemize}
              \item Deployment frequency: Weekly or bi-weekly
          \end{itemize}
    \item \textbf{Interaction Triggers:} Scheduled (manual trigger)
    \item \textbf{Automation Level:} Semi-automated (manual trigger, automated processes)
    \item \textbf{Criticality:} High for system stability and performance
\end{itemize}

% LGTM

\subsection{Other External Interdependencies}

While the Human Benchmark system has a relatively self-contained architecture, there are subtle dependencies worth noting:

\begin{itemize}
    \item \textbf{Browser Feature Dependencies:} The system relies on modern browser capabilities for precise timing measurements. Older browsers may not support the required APIs, potentially affecting reaction time test accuracy.

    \item \textbf{Network Quality Dependencies:} Multiplayer functionality and real-time features depend on stable network connections. Network latency directly impacts the fairness of competitive gameplay, especially for reaction time tests.

    \item \textbf{Database Replication Dependencies:} Information in database must be stored in a consistent state for leaderboard and score tracking. Backups must be performed regularly to prevent data loss.
\end{itemize}

\subsection{Impact of the System on Its Environment}

% LGTM

\subsubsection{Load on Client Devices}

The Human Benchmark system, particularly during reaction and reflex tests, requires dedicated focus from the client device's processing capabilities. This may impact:

\begin{itemize}
    %\item Battery consumption on mobile devices
    \item Available memory and processing power on client devices
    \item Network bandwidth consumption during multiplayer sessions
\end{itemize}

% LGTM

\subsubsection{Data Storage Requirements}

With the growth of the user base and the number of custom games, the system will generate a significant amount of data. Database growth must be monitored and managed.

% LGTM

\subsubsection{Infrastructure Scaling Requirements}

The system is expected to experiance variable loads, especially during peak usage times. System usage must be monitored and scaled accordingly.

%\subsection{Overall Completeness, Consistency, and Coherence}

% ? LGTM
% ! No clue !
% ! Subsection looks really pointless. These things were mentioned before

%The Human Benchmark system provides a way to improve cognitive and motor skill, with comprehensive coverage of:
%\begin{itemize}
 %   \item User management and authentication
 %   \item Game creation and configuration
 %   \item Single-player and multiplayer experiences
 %   \item Performance measurement and tracking
 %   \item Score comparison and leaderboards
%\end{itemize}

%The context model ensures that all business processes required for a complete cognitive testing and training platform are covered, with appropriate attention to the quality and accuracy of measurement, which is central to the system's value proposition.

\section{Perspective Priorities for the Context View}

The following table outlines the priority of different architectural perspectives for the Context View of Human Benchmark, helping to ensure quality attributes are appropriately addressed:
% 🔴 🟡 🟢
% LGTM



\begin{table}[h!]
    \centering
    \caption{Perspective Priorities for the Context View of Human Benchmark}
    \label{tab:perspective_priorities_context}
    \begin{tabular}{p{0.3\textwidth}p{0.15\textwidth}p{0.45\textwidth}}
        % \toprule
        \textbf{Architectural Perspective} & \textbf{Priority} & \textbf{Why needed for Human Benchmark within this View}                                                                                                                                                 \\
        % \midrule
        Accessibility                      & \yellowcircle Medium            & The Context View must identify user interfaces where accessibility is a concern, particularly as the system aims to be worldwide used. \\

        Availability and Resilience        & \redcircle High              & Critical to identify external dependencies that could affect system availability, especially for multiplayer functionality and competition features which require consistent uptime.                    \\

        Development Resource               & \greencircle Low               & While resource constraints exist, they have minimal impact on the Context View's representation of external interfaces.                                                                                 \\

        Evolution                          & \yellowcircle Medium            & The context boundaries should anticipate future expansion of game types and potential integrations with external systems.                                                                               \\

       Internationalization               & \yellowcircle Medium            & The system's context should account for global accessibility, though detailed internationalization concerns belong in other views.                                                                      \\

        Location                           & \yellowcircle Medium            & Geographic distribution of users affects latency considerations for reaction time tests, which must be addressed in the context model.                                                                  \\

        Performance and Scalability        & \redcircle High              & Critical for a reaction-time measurement system; context boundaries must account for performance bottlenecks at interfaces and indicate scaling approaches.                                             \\

        Regulation                         & \greencircle Low               & Limited regulatory concerns for a gaming/cognitive assessment platform in the current scope.                                                                                                            \\

        Security                           & \redcircle High              & The context boundaries must clearly define trust boundaries and authentication/authorization requirements for protecting user data and game integrity.                                                  \\
 
        Usability                          & \yellowcircle Medium            & While important for the system overall, detailed usability concerns are primarily addressed in other views beyond high-level interface identification in the context model.                             \\
        % \bottomrule
    \end{tabular}
\end{table}

\section{Impact of Key Perspectives on the Context View}

\subsection{Availability and Resilience Perspective}

Applying the Availability and Resilience perspective to the Context View of Human Benchmark reveals several critical considerations:

% ? LGTM
% ! No clue ! check if it matches with stuff said previously...

\begin{itemize}
    \item \textbf{Interface Resilience:} The context boundary clearly
          distinguishes which interfaces require high availability (game sessions,
          scores).
          % ! versus those that can tolerate occasional downtime (e.g., leaderboard updates).

    \item \textbf{Dependency Management:} The context model explicitly
          identifies external dependencies that could affect system availability, such
          as the database system and hosting infrastructure
          % ! , with notations about their expected availability characteristics.

          % ! \item \textbf{Graceful Degradation:} The interface definitions account for scenarios where parts of the system might be unavailable, indicating which features could continue functioning with limited capabilities and which require full connectivity.

    \item \textbf{Client-Side Resilience:} The user interface boundary acknowledges potential client-side connectivity issues, with notes about how the system can maintain gameplay integrity despite temporary connection losses.
\end{itemize}

These considerations highlight the important context boundaries to ensure
stakeholders understand availability requirements and dependencies, which is
very important for a system focused on performance measurement and
competitive playing.

\subsection{Performance and Scalability Perspective}

% LGTM

The Performance and Scalability perspective also influences the Context View for
Human Benchmark:

\begin{itemize}
    \item \textbf{Interface Performance Requirements:} The context model
    explicitly identifies interfaces with high performance requirements,
    particularly the user interfaces for reaction time tests where millisecond
    precision is highly important.

    % ! \item \textbf{Scaling Boundaries:} The view clearly shows which components need to scale independently (database or application server especialy during viral surges) versus components that can scale together, informing the high-level system structure.

    % ! \item \textbf{Load Distribution:} The context boundaries account for potential geographic distribution of users and the corresponding need for distributed infrastructure, particularly for minimizing latency in multiplayer scenarios.

    \item \textbf{Resource Consumption Patterns:} The system-environment
    boundary identifies expected resource consumption patterns during peak
    hours, providing important context for infrastructure planning.
\end{itemize}

These considerations are important for a cognitive testing platform where
performance directly impacts the core functionality of our web application.

\subsection{Security Perspective}

The Security perspective shapes the Context View in several important ways:

% LGTM

\begin{itemize}
    \item \textbf{Trust Boundaries:} The context model explicitly identifies
    trust boundaries between the system and external entities, clarifying where
    authentication and authorization controls are required.

    \item \textbf{Data Protection Scope:} The model shows which external
    interfaces handle sensitive user data, requiring encryption or other
    security measures.

    % ! \item \textbf{Authentication Interfaces:} The view specifically highlights the authentication mechanisms and flows at the system boundary, a critical concern for protecting user accounts and score integrity.

    % ! \item \textbf{Tampering Prevention:} The context boundaries acknowledge potential attack vectors for score manipulation or performance enhancement cheating, informing the need for security controls at specific interfaces.
\end{itemize}

For a our platform it is especially important that our users play in a fair
environment and that their accounts are safe. It is important for us that our
system stays safe.

\section{Importance and Justification of the Context View for Human Benchmark}

The Context View is important for the Human Benchmark system for several reasons:

% LGTM

\subsection{Clear System Boundaries}

Human Benchmark's purpose is to provide accurate measurement and comparison of
cognitive and motor skills. The Context View mainly defines the scope of
what the system does and what it does not do, claryfies which external factors might influence
its performance and establishes the boundary between the system's control and external variables.

% \begin{itemize}
%     \item Defining the scope of measurement responsibility (what the system can and cannot accurately measure)
%     \item Clarifying which external factors might influence measurement accuracy (browser differences, network latency)
%     \item Establishing the boundary between the system's control and external variables
% \end{itemize}

% ! This clarity is particularly important for a cognitive testing platform where scientific accuracy depends on understanding system limitations.

% LGTM

\subsection{Interface Requirements Critical to Core Functionality}

The system's functionality relies on the interface characteristics. 
Mainly the precise timing of the interfaces (websockets) for accurate reaction
time measurement and the integrity of score recording interfaces. 

% \begin{itemize}
%     \item Precise timing interfaces are required at the millisecond level for
%     accurate reaction time measurement
%     \item Real-time multiplayer interfaces must ensure fairness in competitive
%     scenarios
%     \item Score recording interfaces must maintain integrity and prevent
%     manipulation
% \end{itemize}

By focusing stakeholder attention on these boundary concerns early, the Context
View ensures the architectural foundations support the system's fundamental
purpose.

% LGTM gud enough

\subsection{External Dependency Management}

The success of Human Benchmark depends on managing several critical external dependencies.
It needs to be ensurerd that web app runs smoothly on various
browsers, different network conditions, and that our databases performance is 
adequate for leaderboard functionality. It also has to be ensured
that systems infrastructure can scale to handle the load during popularity spikes.


The Context View makes these dependencies explicit, allowing stakeholders to assess risks and plan mitigations early in the development process.



\subsection{Stakeholder Alignment}

The Human Benchmark system serves diverse stakeholders including players, 
game creators, and administrators. The Context View provides a common, 
easily understood reference point that helps align stakeholder expectations by:

\begin{itemize}
    \item Clearly communicating system scope and exclusions
    \item Establishing a shared vocabulary for system interactions
    \item Identifying key external interfaces that different stakeholders will interact with
\end{itemize}

This alignment is especially valuable given the system's dual nature as both a measurement tool
 and a competitive gaming platform, where different stakeholders may have different primary interests.

\subsection{Foundation for Further Architectural Work}

As noted in the software architecture literature, 
the Context View serves as a foundation for more detailed architectural views.
 For Human Benchmark, this is particularly important as it:

\begin{itemize}
    \item \textbf{Functional View}: Informs it by identifying required external interfaces
    \item \textbf{Deployment View}: Guides it by clarifying infrastructure dependencies
    \item \textbf{Information View}: Shapes it by defining external data flows
\end{itemize}

Without a clear Context View, subsequent architectural decisions would lack the necessary 
grounding in system boundaries and external constraints.

\section{Cross-View Consistency Considerations}

The Context View of the Human Benchmark system must maintain consistency with other architectural views 
to ensure a coherent overall architectural description. Key considerations include:

% made this into diagram
% \begin{itemize}
%     \item \textbf{Functional View Consistency:} The real-time multiplayer functionality identified in the Functional View must correspond to the WebSocket interface defined in the Context View.

%     \item \textbf{Information View Consistency:} The data flows crossing system boundaries in the Context View must be reflected in the Information View's data models and flows. For instance, the user profile information mentioned in the Context View should be consistently modeled in the Information View.

%     \item \textbf{Deployment View Consistency:} The external systems identified in the Context View (such as the PostgreSQL database and Docker containerization) must be consistently represented in the Deployment View's physical architecture.

%     \item \textbf{Development View Consistency:} The interface technologies mentioned in the Context View (such as WebSockets for real-time communication) should be consistently reflected in the module organization and dependencies described in the Development View.
% \end{itemize}
\begin{figure}[h!]
    \centering
    \includegraphics[width=0.7\textwidth]{AutomationAndSections/Sections/1contextView/Images/CrossView_ContextView.png}
    \caption{Cross-View Consistency Considerations for Context View}
    \label{fig:context_view_crossview}
\end{figure}


% no shit sherlock
%Maintaining this cross-view consistency is essential for ensuring that the architectural description as a whole provides a coherent representation of the Human Benchmark system. 
%Each view provides a different perspective on the same underlying architecture, and inconsistencies between views would indicate potential architectural flaws or documentation issues.

%This document represents one part of a comprehensive Architectural Description for the Human Benchmark system. It should be read in conjunction with other architectural views to gain a complete understanding of the system architecture.

% \end{document}